\subsection{An embedding-dependent uniform bound for ray diagonals}

The preceding counterexamples vary the embedding while retaining its rank.
The following abstract sufficient theorem records what fixed multigraded
algebra data do control. It makes no claim that an arbitrary proposed
branching algebra already satisfies all the stated hypotheses.

\begin{proposition}[P-F16: fixed positive multigrading gives uniform ray bounds]
\label{prop:fixedmultigrading}
Let
\[
 R=\bigoplus_{\eta\in\mathbb Z^s}R_\eta
\]
be a finitely generated complex domain, generated as a complex algebra by
nonzero homogeneous elements $x_1,\ldots,x_m$, with $m\ge1$ and
$\deg(x_i)=\beta_i\in\mathbb Z^s$.
Assume that a real linear functional $\ell$ satisfies
$\ell(\beta_i)>0$ for every $i$. Write
\[
 M=(\beta_1\ \cdots\ \beta_m),\qquad
 \Delta=\max\bigl(\{1\}\cup
 \{|\det M_{I,J}|:\ |I|=|J|\ge1,\ \det M_{I,J}\ne0\}\bigr),
\]
and set
\[
 D=m\Delta,\qquad L=\operatorname{lcm}(1,2,\ldots,D),
 \qquad d=\dim R.
 \tag{P-F16a}
\]
For every $\gamma\in\mathbb Z^s$ with $\ell(\gamma)>0$, form the ray
diagonal, graded by $n$,
\[
 A_\gamma=\bigoplus_{n\ge0} R_{n\gamma}.
\]
Then $A_\gamma$ is generated by homogeneous elements of ray degree at
most $D$, uniformly in $\gamma$.

Suppose further, for the rays under consideration, that $A_\gamma\ne\mathbb C$,
that $A_\gamma$ is Cohen--Macaulay, and that
$f_\gamma(n)=\dim_\mathbb C R_{n\gamma}$ is a global quasi-polynomial
on all $n\ge0$. If $r_\gamma=\dim A_\gamma$, then $A_\gamma$ has an
h.s.o.p. consisting of $r_\gamma$ elements all of ray degree $L$.
Its Hilbert series therefore has a positive form
\[
 \sum_{n\ge0}f_\gamma(n)T^n
 =\frac{H_\gamma(T)}{(1-T^L)^{r_\gamma}},
 \quad H_\gamma\in\mathbb Z_{\ge0}[T]\setminus\{0\},
 \quad \deg H_\gamma<r_\gamma L.
 \tag{P-F16b}
\]
Consequently its positivity and saturation indices satisfy
\[
 s(f_\gamma)\le p(f_\gamma)
 \le\max\{0,\deg H_\gamma-L\}
 \le(r_\gamma-1)L\le(d-1)L.
 \tag{P-F16c}
\]
The last bound depends on the fixed presentation of $R$ and its weight
matrix, and is independent of $\gamma$. It is an existence bound;
it is not asserted to be polynomial in the binary encoding of $M$.
\end{proposition}

\begin{proof}
The positivity of $\ell$ first gives $R_0=\mathbb C$ and ensures that
every multihomogeneous component of the polynomial algebra
$P=\mathbb C[X_1,\ldots,X_m]$, with $\deg X_i=\beta_i$, is finite
dimensional. Indeed, if $M u=\eta$ with $u_i\ge0$, then
\[
 \sum_i u_i\ell(\beta_i)=\ell(\eta),
\]
which bounds every $u_i$ when such a solution exists. The homogeneous
surjection $P\longrightarrow R$, $X_i\mapsto x_i$, restricts to a
surjection of ray diagonals: a preimage of a multihomogeneous element
can be replaced by its component of the same multidegree. Thus it
suffices first to bound the generating degrees of the semigroup
\[
 S_\gamma=\{(u,n)\in\mathbb Z_{\ge0}^{m+1}:Mu=n\gamma\}.
 \tag{P-F16d}
\]
If this semigroup has no element with $n>0$, its algebra and $A_\gamma$
are just $\mathbb C$, and the first assertion is immediate. Otherwise
the polytope
\[
 Q_\gamma=\{u\in\mathbb R_{\ge0}^m:Mu=\gamma\}
\]
is nonempty. It is bounded, since
$u_i\le\ell(\gamma)/\ell(\beta_i)$.

We give the required rational-polyhedral argument explicitly.
At any vertex $v$ of $Q_\gamma$, the columns of $M$ indexed by
$J=\{i:v_i>0\}$ are linearly independent. Otherwise a nonzero
vector supported on $J$ in the kernel of $M$ would permit sufficiently
small perturbations of $v$ in both directions, contradicting extremality.
If $q=|J|$, select $q$ rows $I$ for which $M_{I,J}$ is nonsingular.
Cramer's rule shows that $|\det M_{I,J}|v$ is integral. Hence the
primitive integral generator of the ray through $(v,1)$ has last
coordinate at most $\Delta$.

For clarity, every point of $Q_\gamma$ is a convex combination of
these vertices without an additional integrality assumption. If its
positive-support columns are dependent, choose a nonzero supported
kernel vector. Such a vector has both positive and negative entries,
since applying $\ell$ to $M u=0$ forbids a nonzero nonnegative or
nonpositive kernel vector. Move in its two directions until a positive
coordinate first reaches zero. The original point is a convex
combination of these two endpoints, each with smaller support.
Induction on support size terminates at points with independent
support columns. Such a point is a vertex: in any convex decomposition,
its zero coordinates remain zero, and independence makes both endpoints
equal to it. There are only finitely many possible independent supports,
so there are finitely many such vertices.

Consider now the real cone
\[
 C_\gamma=\{(u,n)\in\mathbb R_{\ge0}^{m+1}:Mu=n\gamma\}.
\]
Its only point with $n=0$ is zero, by positivity of $\ell$.
Dividing a nonzero point by $n$ and using the preceding vertex
decomposition shows that $C_\gamma$ is generated by the finitely many
primitive integral rays just found. Its dimension $e$ is at most $m$:
the last coordinate is determined by
$n=\ell(Mu)/\ell(\gamma)$, so projection to $u$ is injective.
In particular the positivity assumptions exclude the rank-zero matrix
case.

Every nonzero $w\in S_\gamma$ can be expressed as a nonnegative real
combination of at most $e$ linearly independent primitive ray vectors,
say
\[
 w=\sum_{j=1}^k a_jv_j,\qquad a_j\ge0,\qquad k\le e\le m.
\]
To justify independence, begin with any finite ray expression.
If the active rays are dependent, choose a dependence, which has both
signs because their last coordinates are positive. Subtract a suitable
multiple of it from the coefficients until one active coefficient
becomes zero, retaining nonnegativity. Repeating yields independent
active rays.
Subtract the integer multiples $\lfloor a_j\rfloor v_j$ and obtain
the integral residual
\[
 w_0=\sum_{j=1}^k\{a_j\}v_j\in S_\gamma,
 \qquad n(w_0)<k\Delta\le m\Delta=D.
 \tag{P-F16e}
\]
This is the fundamental-parallelepiped argument, with its degree
estimate made explicit. The primitive vectors themselves have ray
degree at most $\Delta\le D$. Hence $S_\gamma$ is generated by its
elements of degrees at most $D$. There are only finitely many such
elements, by the same positive-functional bounds on their coordinates.
Its semigroup algebra is the polynomial ray diagonal; passing to its
quotient $A_\gamma$ proves the uniform generating-degree assertion.

Fix now a ray satisfying the additional assumptions, and abbreviate
$A=A_\gamma$ and $r=\dim A$. Choose nonzero homogeneous algebra
generators $y_1,\ldots,y_t$ of degrees $1\le b_j\le D$.
Let $B=\mathbb C[A_L]$, the subalgebra generated by its degree-$L$
component. This component is finite dimensional. Since $L$ is divisible
by every $b_j$, the equation
\[
 Y^{L/b_j}-y_j^{L/b_j}=0
\]
is a monic equation for $y_j$ with coefficients in $B$.
Thus $A$ is integral and finite over $B$: adjoining finitely many
integral elements gives a finite module, as is seen by reducing powers
using their monic equations. The domain $B$ has dimension $r$ and is
standard graded when its degree $L$ is renamed degree one.

Here is a direct homogeneous Noether-normalization argument giving
parameters of that same degree. If a standard graded domain $B$ is
generated by $z_1,\ldots,z_q$ of degree one and $q>\dim B$, choose a
nonzero homogeneous polynomial relation $F(z_1,\ldots,z_q)=0$.
Over the infinite field $\mathbb C$, choose $c_1,\ldots,c_{q-1}$ with
$F(c_1,\ldots,c_{q-1},1)\ne0$. Such a choice exists because
dehomogenizing a nonzero homogeneous polynomial at its last variable
gives a nonzero polynomial. Put $z_i'=z_i-c_i z_q$ for $i<q$.
After substituting $z_i'+c_i z_q$ in $F$, its leading coefficient
as a polynomial in $z_q$ is the displayed nonzero scalar.
Rescaling makes this a monic relation for $z_q$ over
$B'=\mathbb C[z_1',\ldots,z_{q-1}']$. Therefore $B$ is finite over
$B'$, while $\dim B'=\dim B$. Induction on the number of generators
terminates with $r$ algebraically independent linear forms
$\theta_1,\ldots,\theta_r$ such that $B$, and hence $A$, is finite
over $\mathbb C[\theta_1,\ldots,\theta_r]$. These forms all have
original ray degree $L$. They are an h.s.o.p. for $A$.

We spell out the Cohen--Macaulay step to keep its hypothesis visible.
Put $\mathfrak m=A_{>0}$. The quotient by all the $\theta_i$ is
finite dimensional, because $A$ is finite over their polynomial
algebra. It is connected and nonnegatively graded; every element of
positive degree is nilpotent. Thus $\mathfrak m$ is the unique maximal
ideal above $(\theta_1,\ldots,\theta_r)$, and the local ring
$A_{\mathfrak m}$ has dimension $r$ (equivalently, lift a maximal
prime chain from the parameter polynomial algebra by going up).
The Cohen--Macaulay regular-parameter criterion therefore says that
$\theta_1,\ldots,\theta_r$ is a regular sequence in $A_{\mathfrak m}$;
an exact primary reference is the Stacks Project,
Lemma 10.104.2, Tag \texttt{02JN},
\url{https://stacks.math.columbia.edu/tag/02JN}.
It is also a regular sequence in $A$. Indeed the kernels of the
successive homogeneous multiplications are graded. A nonzero
homogeneous element of such a kernel cannot disappear on localization
at $\mathfrak m$: multiplying it by an element outside $\mathfrak m$
has a nonzero lowest-degree term, coming from the nonzero scalar
constant coefficient of that element. Local injectivity therefore
implies graded global injectivity.

The exact sequences for multiplication by these regular homogeneous
parameters give
\[
 \operatorname{Hilb}_{A/(\theta_1,\ldots,\theta_r)}(T)
 =(1-T^L)^r\operatorname{Hilb}_{A}(T).
\]
The quotient is finite dimensional and has nonnegative integer graded
dimensions and constant coefficient one. Its Hilbert series, which
is a polynomial, is the required $H_\gamma(T)$. The global quasi-polynomial assumption
and the properness argument of P-F9 give
$\deg H_\gamma<rL$.
In this equal-denominator-degree situation a direct argument improves
the general bound of P-F11. Write $h=\deg H_\gamma$ and
$H_\gamma(T)=\sum_{i=0}^h H_iT^i$.
For each residue $b$ modulo $L$, the corresponding eventual constituent
of $f_\gamma$ is
\[
 q_b(X)=\sum_{\substack{0\le i\le h\\i\equiv b\pmod L}}
 H_i\binom{(X-i)/L+r-1}{r-1}.
\]
This follows by expanding $(1-T^L)^{-r}$ at sufficiently large
indices in that residue. Since $f_\gamma$ is globally
quasi-polynomial, these are its constituents everywhere: any
constituent for a common refinement of its period with $L$ agrees
with the displayed polynomial at infinitely many integers and hence
is identical to it. Put $c=\max\{0,h-L\}$. After the integer shift
$c$, each binomial polynomial appearing in $q_b(X+c)$ is
\[
 \frac{1}{(r-1)!}\prod_{j=1}^{r-1}
 \left(\frac{X}{L}+\frac{c-i}{L}+j\right).
\]
All its coefficients are nonnegative, because
$c-i+L\ge c-h+L\ge0$ and $j\ge1$.
For $r=1$ the empty product is one. Nonnegative linear combinations
retain that property, and shifting merely permutes the residue
labels. Therefore $p(f_\gamma)\le\max\{0,h-L\}\le(r-1)L$,
using $h<rL$; when $r=1$ this bound is zero.
Since $A$ is a finitely generated subdomain of $R$,
$r=\operatorname{trdeg}_\mathbb C\operatorname{Frac}(A)
\le\operatorname{trdeg}_\mathbb C\operatorname{Frac}(R)=d$.
Finally, a nonzero polynomial with nonnegative coefficients is
strictly positive at every positive integer, so coefficient positivity
after a shift implies the source's strict saturation after the same
shift. This proves all assertions.
\end{proof}

\paragraph{Precise use and remaining gate.}
For a fixed embedding, an application requires a single finitely
generated, positively multigraded domain whose $(\lambda,\pi)$
components are the actual branching multiplicity spaces, together with
Cohen--Macaulayness and global quasi-polynomiality for the ray diagonals.
A polynomial branching algebra is a possible input to verify.
The full algebra $\mathbb C[GL_m/U]$ has determinant units and cannot
be inserted into the positivity hypothesis on $\ell$ without addressing
its grading. No such verification is supplied by the abstract theorem.

The quantitative data required here are the generator number, their
multidegrees, the resulting minors, and the parameter degree $L$.
They can change when the embedding changes, as in P-F14--P-F15.
The formula $L=\operatorname{lcm}(1,\ldots,m\Delta)$ is a finite
explicit bound; even a polynomial numerical bound on $m\Delta$
does not by itself turn this particular formula into a polynomial
bound for $L$. A uniform PH3 or PH2 conclusion needs stronger control
of parameter degrees, or another bound avoiding this least-common-multiple
growth. The theorem supplies a corrected sufficient gate and a fixed-data
uniformity statement, not the missing uniform algorithm or a claim of
historical priority.
